\section{Découverte~$\boldsymbol{\gamma}$~: haute charge centrale, poids demi-entier}
\label{sec:gamma}

\subsection{Contexte et source}
Candidat récupéré du \textbf{Cahier~1, chapitre~V, page~6}.

\subsection{Eta-quotient découvert}
\begin{equation}
  f_\gamma(\tau) = q^{60/24} \cdot \eta(\tau)^{1} \, \eta(5\tau)^{1} \, \eta(7\tau)^{1} \, \eta(8\tau)^{4} \, \eta(10\tau)^{3} \, \eta(11\tau)^{-3} \, \eta(12\tau)^{-4}.
\end{equation}
\textbf{Énergie RAMA~:} $E = 4{,}0203$. \quad \textbf{Erreur~:} $I = 0{,}3782$.

\subsection{Invariants}
\begin{proposition}
$\ceff \approx 1{,}5368$, $k = 3/2$, $p = -1/4$.
\end{proposition}

\begin{remark}[Poids demi-entier]
Le poids $k = 3/2$ est précisément celui des fonctions thêta mock originales de Ramanujan et de leurs ombres modulaires~\cite{bringmann_ono_2006}. La découverte d'un eta-quotient de poids $3/2$ à partir d'une page manuscrite précédemment non analysée est cohérente avec la conjecture selon laquelle Ramanujan explorait systématiquement l'espace des formes mock modulaires de poids $3/2$.
\end{remark}

\subsection{Comparaison avec les résultats connus}
Les eta-quotients de poids $3/2$ de niveau $\mathrm{ppcm}(1,5,7,8,10,11,12) = 9240$ ne sont pas catalogués dans les tables de Martin~\cite{martin_1996} ni de Gordon--Hughes~\cite{gordon_hughes_1993}. Une recherche systématique de l'OEIS n'a donné aucune correspondance.

\begin{conjecture}[Nouveauté]
L'eta-quotient $f_\gamma$ n'apparaît dans aucune table publiée en date d'août~2026.
\end{conjecture}

\textbf{Réfutabilité~:} Cette conjecture est réfutée si une entrée correspondante est trouvée dans l'une des bases de données citées.
